\documentclass[11pt]{article}
\usepackage[T1]{fontenc}
\usepackage[margin=1in]{geometry}
\usepackage{amsmath,amssymb}
\usepackage{booktabs,tabularx,longtable}
\usepackage{xurl}
\usepackage[hidelinks]{hyperref}
\newcolumntype{L}{>{\raggedright\arraybackslash}X}
\setlength{\emergencystretch}{3em}
\newcommand{\code}[1]{\texttt{\detokenize{#1}}}
\title{CorrPFN: Selecting Useful Sample--Formula Pairs\\
\large Design, Implementation, and First Recipe Results}
\author{}
\date{14 November 2026}

\begin{document}
\maketitle

\begin{abstract}
The new experiment asks a specific question: can informed training become more
useful if it focuses on real samples where a particular target formula predicts
better than a generic model? Instead of choosing formulas without considering
the source sample, we keep successful real-sample--formula pairs. A synthetic
example is generated by choosing one retained pair, perturbing its parent's
composition, keeping its experimental conditions, and applying its assigned
formula. The first recipe tests ten formulas on 760 real EPiT samples against
five-fold query predictions from the generic step-10,000 checkpoint. It retains
2,300 of 7,600 pairs, giving a training mixture of approximately 30.26\% informed
and 69.74\% generic. This document describes the implemented recipe and its
checks. It does not claim that the new trained model has improved: a completed
training run and a model comparison are not established by these results.
\end{abstract}

\tableofcontents
\newpage

\section{Why the project changed}
The earlier work examined where informed models, especially the model trained
with an informed proportion of 1.0, failed to improve over generic training.
The aim was to find specific compositions or test conditions that were poorly
represented, or poorly described by the target formulas.

A later diagnostic applied the saved formulas to real features. Each real
sample received one formula, with assignments following the old training
probabilities. Repeating this assignment 1,000 times measured the combined
formula mixture, rather than choosing the best formula after seeing each answer.

The five-fold comparison found that the combined formulas had lower average
absolute error than generic on 237 of 760 samples, lower error than pretrained
on 148, and lower error than both on 120. These counts used an improvement of
more than 1 mV. Among the 120 samples where formulas beat both models, informed
training improved over generic on 56, worsened on 60, and was within 1 mV on four.
The generic and informed errors in that diagnostic were averaged over
checkpoints 500--5500; pretrained was a separate fixed baseline.

These observations motivated the new experiment. They suggest that some formula
predictions may be useful even where informed training has not delivered their
benefit. They do not prove that the old target-mixing weight caused the failures,
or that every poorly predicted synthetic example harms training.

\textbf{The new recipe is a separate comparison.} It refits coefficients on all
760 real samples, evaluates every sample--formula pair, and uses the single
generic step-10,000 checkpoint. The earlier count of 120 is therefore not the
number of parents or pairs in the new recipe.

\section{The agreed design}
The decisions for the first experiment are deliberately simple:
\begin{enumerate}
\item Fit the ten formula families on all 760 real samples.
\item Obtain one generic query prediction for each sample through five-fold CV.
\item Keep a sample--formula pair if its absolute error is strictly lower than
      the generic model's error on that sample.
\item Give every retained pair equal sampling probability.
\item Use one real parent per synthetic example. Keep its composition and
      experimental conditions together.
\item Perturb composition using the existing multiplicative-noise method.
\item Apply the selected formula with fixed fitted coefficients. Do not mix its
      target with an SCM target.
\item Set the informed training proportion to the fraction of pairs retained.
\end{enumerate}

Here, an SCM is the random structural causal model used by the generic synthetic
generator. Generic training continues to use its usual generator. Removing SCM
mixing concerns the targets of informed examples only.

\section{Real data and generic predictions}
\subsection{Inputs and target}
The recipe uses the reviewed EPiT export with 760 usable real observations.
Each model input has 70 physical feature columns: 24 composition columns and
46 columns describing environment, processing, and experimental history.
The measured target is pitting potential in mV relative to SCE.

The preparation script reads the existing reviewed features and target loader.
It does not redo source preprocessing, invent missing measurements, or change
which 760 rows are accepted. Numeric threshold targets retain the existing
loader convention; they are not newly treated as censored observations in this
experiment.

\subsection{Five-fold query predictions}
The generic reference is
\begin{quote}\small
\path{checkpoints/tabicl_s1_regression_baseline/step-10000.ckpt}.
\end{quote}
The existing exploratory five-fold layout is reused. Each fold has 152 queries
and the remaining 608 samples as labelled context. Fold 1 keeps the historical
outer query set. The original development rows are split into the other four
query folds using seed 42. These are the established exploratory partitions,
not newly constructed publication-held-out folds.

Each real sample appears as a query exactly once. It appears in context in the
other four folds, but those appearances do not provide additional query
predictions. Thus, \textbf{no averaging of four generic predictions is needed}.
For selection, the five query exports are joined into one list of 760 predictions.

Inference uses the existing native TabICL input path, eight estimators, seed 42,
median regression output, and CPU execution. No pretrained-model prediction is
used in the new pair-retention decision.

\section{Fitting and applying the formulas}
\subsection{Formula inventory}
The experiment keeps the nine promoted base rules and the heat-category
extension already used in the preceding EPiT work. It changes their fitted
coefficients and the way successful pairs are selected; it does not add new
physical formula families.

\begin{table}[htbp]
\centering\small
\begin{tabularx}{\linewidth}{lL}
\toprule
Rule identifier & Main ingredients \\
\midrule
\code{pren_n_linear} & Cr--Mo--W--N material score and environment terms. \\
\code{cr_mow_n_synergy} & Adds a chromium--molybdenum/tungsten interaction. \\
\code{threshold_saturation} & Chromium threshold and saturation behaviour. \\
\code{pren_n_improved_environment} & Separate chloride, temperature, and acidic-pH terms. \\
\code{mo_n_acid_repassivation} & Adds an Mo/W--N interaction under acidic conditions. \\
\code{mns_inclusion_penalty} & Adds an Mn--S inclusion-related penalty. \\
\code{coupled_breakdown} & Describes environmental attack relative to material protection. \\
\code{fe_ni_cr_threshold} & Different chromium thresholds for Fe-like and Ni-like compositions. \\
\code{method_aware_pren_n} & Includes a correction for the reported test method. \\
\code{extend_heat_categories} & Extends the method-aware base with heat/cooling categories. \\
\bottomrule
\end{tabularx}
\caption{The ten existing formula families used in the selected-pair recipe.}
\end{table}

\subsection{All-row fitting is intentional}
The new fitting script explicitly supplies all 760 rows to the fitting routines.
It does not apply the old Fe/Ni-only eligibility filter before fitting. This
means the equations are also fitted and tested on aluminium, HEA, and other
alloys, although their original physical assumptions do not automatically become
valid for every alloy family.

For each of the nine base rules, feature imputation and scaling are fitted on
these rows and saved with the coefficients. The fitting objective uses
standardized target ranks. In plain language, it first fits how samples should
be ordered from lower to higher EPiT. Base-rule coefficients are nonnegative,
sum to one, and obey the existing upper bounds. The existing penalty toward the
rule's reference coefficients remains 0.1.

The heat extension uses the same all-row data and its existing fitting method.
Its penalty remains 0.3; a new penalty search is not performed. Its extension
coefficients are not assumed to obey the base rules' sum-to-one constraint.
All fitted states are saved separately from the old development-only recipes.

\subsection{Converting formula scores to mV}
A formula first returns a score, not necessarily a value in mV. A separate
positive affine conversion is fitted for each formula:
\begin{equation}
 \widehat y_{ir}=a_r+b_r s_r(x_i),\qquad b_r\geq 0,
\end{equation}
where $x_i$ contains the real features and $s_r$ is rule $r$'s score. The
conversion uses all 760 measured targets in this experiment. Its intercept and
slope are saved with the formula.

This conversion matters: comparing an uncalibrated rank score directly with a
measurement in mV would not give a meaningful absolute error. During synthetic
generation, both the fitted formula and its mV conversion stay fixed.

\section{Selecting the useful pairs}
Let $y_i$ be sample $i$'s measured target, $g_i$ its generic query prediction,
and $\widehat y_{ir}$ the prediction from formula $r$. Define
\begin{align}
 e_i^{\mathrm{generic}} &= |g_i-y_i|,\\
 e_{ir}^{\mathrm{formula}} &= |\widehat y_{ir}-y_i|.
\end{align}
The retained set is
\begin{equation}
 \mathcal{S}=\{(i,r):e_{ir}^{\mathrm{formula}}<e_i^{\mathrm{generic}}\}.
\end{equation}

There is no extra 1 mV or 25 mV threshold in this new selection rule. Any strictly
smaller error qualifies; equal errors do not. This follows the agreed first
experiment rather than adding a new selection threshold.

A pair is a specific real row and a specific formula. A formula can qualify for
one sample and fail for another. Likewise, one sample can retain several rules.

\subsection{Observed selection results}
\begin{table}[htbp]
\centering
\begin{tabular}{lr}
\toprule
Quantity & Result \\
\midrule
Real samples & 760 \\
Formulas per sample & 10 \\
Tested pairs & 7,600 \\
Retained pairs & 2,300 \\
Real samples with at least one retained rule & 424 \\
Real samples with no retained rule & 336 \\
Retained pairs improving error by more than 25 mV & 1,810 \\
\bottomrule
\end{tabular}
\caption{Results from the saved all-row recipe, not from a newly trained model.}
\end{table}

\begin{table}[htbp]
\centering\small
\begin{tabular}{lr}
\toprule
Formula & Retained pairs \\
\midrule
\code{pren_n_linear} & 221 \\
\code{cr_mow_n_synergy} & 219 \\
\code{threshold_saturation} & 235 \\
\code{pren_n_improved_environment} & 219 \\
\code{mo_n_acid_repassivation} & 229 \\
\code{mns_inclusion_penalty} & 209 \\
\code{coupled_breakdown} & 225 \\
\code{fe_ni_cr_threshold} & 239 \\
\code{method_aware_pren_n} & 230 \\
\code{extend_heat_categories} & 274 \\
\midrule
Total & 2,300 \\
\bottomrule
\end{tabular}
\caption{All ten rules retain at least some pairs.}
\end{table}

\begin{table}[htbp]
\centering
\begin{tabular}{lrr}
\toprule
Material family & Eligible real parents & Retained pairs \\
\midrule
Fe alloys & 322 & 1,605 \\
NiCrMo alloys & 41 & 249 \\
Al alloys & 28 & 220 \\
HEA & 15 & 95 \\
Other & 18 & 131 \\
\midrule
Total & 424 & 2,300 \\
\bottomrule
\end{tabular}
\caption{Parents available to the new informed sampler. ``Eligible'' here means
that at least one rule beat generic, not that the source measurement is independently verified.}
\end{table}

\section{What the old sampler actually did}
Inspecting the implementation corrected an earlier assumption: the old EPiT
sampler did \textbf{not} average two compositions 50/50.
It selected one composition template, perturbed it, and independently selected
an experimental-condition record. It copied that record's environment and
history fields together. As a result, composition could come from sample A
while conditions came from sample B.

The old informed feature profile was restricted to Fe and NiCrMo sources.
Formula choice was made at the synthetic-table level. These details explain why
the new project required changes to both source selection and target assignment,
rather than simply filtering the existing composition list.

After this inspection, the agreed first implementation changed to
\textbf{one parent per informed example}. No second-parent search, parent
interpolation, or fallback between two parents is part of the implemented design.

\section{The new informed sampler}
\subsection{Choose a pair, not just a parent}
The sampler draws uniformly from the retained list:
\begin{equation}
 P((i,r))=\frac{1}{|\mathcal{S}|}.
\end{equation}
If real sample $i$ has $k_i$ retained rules, its total parent probability is
\begin{equation}
 P(i)=\frac{k_i}{|\mathcal{S}|}.
\end{equation}
Thus, a parent with three successful rules is sampled three times as often as a
parent with one. There is no separate weighting by the size of the improvement,
by publication, or by alloy family.

\subsection{Keep the parent's conditions}
The selected real row supplies both its composition and its experimental
conditions. Temperature, chloride, pH, heat treatment, surface preparation,
test-method information, and missing values remain associated with that parent.
The conditions are not drawn independently from another row.

Category codes are constructed consistently from the reviewed data. Missing
model-input features remain marked as missing until the existing input path
fills them using context rows. Formula evaluation uses the saved formula-specific
imputation and scaling states; these are not refitted on each synthetic table.

\subsection{Perturb composition}
For a positive parent element amount $c_j$, the existing noise method is
\begin{equation}
 \widetilde c_j=c_j\exp(\epsilon_j),\qquad
 \epsilon_j\sim\mathcal{N}(0,\sigma^2).
\end{equation}
The perturbed composition is closed to 100 wt.\%:
\begin{equation}
 c_j'=100\frac{\widetilde c_j}{\sum_k\widetilde c_k}.
\end{equation}
An element that is zero in the parent stays zero. The launch configuration keeps
$\sigma=0.30$, the preceding v5 default. This is the standard deviation of noise
in log space; it is not a strict $\pm30\%$ bound. A zero-perturbation diagnostic
leaves the original composition unchanged so saved predictions can be replayed.

\subsection{Calculate the assigned formula's target}
The formula selected with the parent is applied to the newly perturbed features.
Its coefficients and mV conversion remain fixed:
\begin{equation}
 y'=a_r+b_r s_r(c',\text{parent conditions}).
\end{equation}
There is no random coefficient variation and no extra SCM contribution to this
target. Different rows in one informed table can have different assigned rules.
The table's completed mV targets receive one shared standardization before
training; each rule is not independently rescaled on its subset of the table.

A formula that works on the original sample is not guaranteed to work on every
perturbation around it. That local-transfer assumption is part of the experiment,
not something pair selection proves.

\section{Choosing the generic/informed proportion}
The agreed data-based starting proportion is
\begin{equation}
 p_{\mathrm{informed}}=
 \frac{\text{retained pairs}}{\text{tested pairs}}
 =\frac{2300}{7600}=0.3026315789.
\end{equation}
The remaining proportion, approximately 69.74\%, uses generic generation.
This rule is a starting design choice. The fraction of successful pairs does
not establish the best training ratio.

\subsection{How this fits the existing training batches}
The existing training code mixes \emph{whole synthetic tables} in groups. Each
table has its own context and query rows. The implementation preserves that
structure; it does not put unrelated generic rows and informed rows into the
same table merely to reach a percentage.

With $G$ table groups in the global batch and indexed batch $t$, the number of informed
groups is
\begin{equation}
 n_t=\lfloor(t+1)G p_{\mathrm{informed}}\rfloor
      -\lfloor tG p_{\mathrm{informed}}\rfloor.
\end{equation}
This carries rounding across batches. It avoids permanently rounding 30.26\%
to a nearby representable share. The requested ratio is approached in table-group
counts over time; it is not an exact fraction of rows in every individual batch.
The global batch is then partitioned across GPUs, preserving its tables and
their order when the GPU count changes.

\section{Implementation and saved artifacts}
\subsection{Preparation}
\begin{quote}\small
\path{scripts/epit_pipeline/prepare_selected_pairs.py}
\end{quote}
This script fits the all-row formulas, exports generic 10k predictions for the
five query folds, calculates pair errors, and saves the retained list.
It checks row coverage, query/context separation, duplicate identities, and
matching measured targets. Completed fold exports can be reused after their
saved hashes and identities are checked.

Its output directory is
\begin{quote}\small
\path{corrosion_datasets/analysis/epit_pipeline/selected_pairs_all760_generic10k/}.
\end{quote}
Important files are:
\begin{description}
\item[\code{fitted_formulas.json}] Formula coefficients, preprocessing states,
fit-row identities, and mV conversions.
\item[\code{formula_predictions.csv}] Every formula's prediction on every real row.
\item[\code{generic_cv_predictions.csv}] One generic query prediction per real row.
\item[\code{all_pairs.csv}] Formula and generic errors for all 7,600 pairs.
\item[\code{retained_pairs.csv}] The 2,300 selected pairs.
\item[\code{summary.json}] Counts and the calculated informed proportion.
\item[\code{manifest.json}] Data, checkpoint, implementation, and fold identities.
\end{description}

\subsection{Sampler and training connection}
\begin{quote}\small
\path{src/tabicl/prior/epit_selected_pairs.py}
\end{quote}
\code{SelectedPairSampler} loads the frozen recipe, checks its inputs, samples
pairs, perturbs compositions, and evaluates the saved formulas. It returns the
parent indices, pair indices, and formula names for inspection.

\code{SelectedPairPriorMixture} connects this sampler to the existing original
generic mixture. It derives the proportion from the pair inventory, uses the
indexed group allocation above, and preserves the generic generation path.
Recipe files and formula implementation hashes are included in the sampling
identity used for checkpoint compatibility.

The opt-in training argument is \code{--epit_selected_pair_dir}, added through
\path{src/tabicl/train/train_config.py} and
\path{src/tabicl/train/run.py}. Without this option, existing samplers remain
available. The old v5 target mixer is not used by this experiment.

\subsection{Training launcher}
\begin{quote}\small
\path{scripts/epit_pipeline/train_selected_pairs.py}
\end{quote}
The launcher reuses the established EPiT architecture, optimizer, and batch
configuration through the old launcher's planning-only path. It then selects
the new sampler, disables coefficient variation, and starts a separate run.
It does not enter the old launcher's automatic model-selection or evaluation
stages.

The default is fresh training, not weight initialization from generic 10k.
The generic 10k checkpoint is the \emph{selection reference}, not the new model's
initial weights. The default launch settings include:
\begin{table}[htbp]
\centering
\begin{tabular}{ll}
\toprule
Setting & Value \\
\midrule
Training steps & 10,000 \\
Device & One CUDA GPU \\
Global batch size & 512 synthetic tables \\
Table group size / micro-batch size & 4 / 4 \\
Maximum table length & 1,024 rows \\
Learning rate & $10^{-4}$ \\
Schedule & Cosine with 2\% warmup \\
Schedule horizon & 10,000 steps, independent of training duration \\
Data-loader workers & 4 \\
Composition perturbation & $\sigma=0.30$ \\
Coefficient variation & 0 \\
Informed target & Assigned formula only \\
Temporary / permanent checkpoint intervals & 100 / 500 steps \\
\bottomrule
\end{tabular}
\caption{Default launch settings. The informed share comes from the saved pairs.}
\end{table}

The command still uses the old \code{v4} dispatch setting to avoid entering the
old v5 mixer. This does \emph{not} remove the heat formula: the selected-pair
sampler loads all ten newly fitted equations from its own artifact.

\section{Checks completed}
Focused checks cover the following behaviour:
\begin{enumerate}
\item With perturbation off, the new evaluator reproduces saved predictions for
      sampled pairs spanning all ten rules.
\item With perturbation on, composition totals remain 100\%, original zeros stay
      zero, and experimental conditions stay with the parent.
\item Repeating the same seed reproduces sampled features and targets.
\item Pair selection is uniform over the retained list, and the group allocator
      carries fractional informed shares across batches.
\item A synthetic training batch can be generated with finite model inputs and
      targets. Resuming its indexed sampling state reproduces the next batch.
\item Generic source tables match the existing generator under matched settings.
      Informed target generation does not call the SCM target-generation method.
\item Pair-retention checks reject duplicate identities, inconsistent targets,
      and nonfinite predictions; equal-error pairs are discarded.
\item The launch plan is checked for CPU and CUDA, fixed coefficients,
      formula-only targets, the selected-pair ratio, and fresh initialization.
\end{enumerate}
Nine focused sampler/selection/stream checks and two launch-plan checks passed.
These are implementation checks, not evidence of improved model accuracy. They
do not constitute a completed full-training or full-repository test run.

\section{Commands and output locations}
Run commands from \path{/home/fcolanto/projects/tabicl}.

\subsection{Prepare the recipe}
This step has already been completed for the reported recipe:
\begin{verbatim}
PYTHONPATH=src:. .venv/bin/python \
  -m scripts.epit_pipeline.prepare_selected_pairs
\end{verbatim}

\subsection{Check the launch without starting training}
\begin{verbatim}
PYTHONPATH=src:. .venv/bin/python \
  -m scripts.epit_pipeline.train_selected_pairs --dry-run
\end{verbatim}

\subsection{Start the separate training run}
\begin{verbatim}
PYTHONPATH=src:. .venv/bin/python \
  -m scripts.epit_pipeline.train_selected_pairs
\end{verbatim}
Checkpoints are written under
\begin{quote}\small
\path{checkpoints/epit_selected_pairs/epit_selected_pairs_all760_generic10k_v1/}.
\end{quote}
The training plan and log are written under
\begin{quote}\small
\path{corrosion_datasets/analysis/epit_pipeline/selected_pair_training/epit_selected_pairs_all760_generic10k_v1/}.
\end{quote}
The launcher refuses to replace a nonempty run directory. A different
\code{--run-name} creates a separate experiment. The launcher does not currently
provide a dedicated resume flag; sampling-state resume has been tested at the
underlying prior level. Automatic post-training evaluation is not included.

\section{What the results do and do not establish}
\subsection{All real targets helped construct this recipe}
The formulas, their mV conversions, and pair selection use all 760 measured
targets. Generic predictions exclude the query sample from context, but that
does not make the formula side held out. The resulting comparison is
intentionally a way to construct a training recipe, not an independent estimate
of formula accuracy.

The historical 152-row outer set is also used in fitting and selection now.
It cannot subsequently be described as an untouched final test for this new
recipe. Testing on the same 760 rows can describe model behaviour, but an
independent test needs separate data or a construction that excludes the test
rows throughout fitting and selection.

\subsection{Remaining scientific questions}
A retained formula may still have a large absolute error if generic is worse.
A formula may also match one measurement by chance. Multiple retained rules for
the same parent can disagree around that parent. No extra accuracy threshold,
publication balancing, or neighbourhood validation has been added to this first
experiment.

The 30.26\% informed proportion is an agreed empirical starting rule, not an
optimized training setting. Fixed coefficients and formula-only targets simplify
the experiment, but their benefit must be measured rather than assumed.

The next result to establish is whether the trained model improves under this
recipe, and on which samples. Until that comparison exists, the supported claim
is that the recipe has been constructed and connected to training---not that
selective informed training has already succeeded.

\section{Repository sources for this document}
The description is based on the local implementation and saved artifacts:
\begin{itemize}
\item \path{scripts/epit_pipeline/prepare_selected_pairs.py}
\item \path{scripts/epit_pipeline/formula_extension_study.py}
\item \path{scripts/epit_pipeline/calibrate_target_rules.py}
\item \path{scripts/epit_pipeline/exploratory_cv.py}
\item \path{src/tabicl/prior/epit_feature_generator.py}
\item \path{src/tabicl/prior/epit_composition_profile.py}
\item \path{src/tabicl/prior/epit_selected_pairs.py}
\item \path{scripts/epit_pipeline/train_selected_pairs.py}
\item \path{scripts/epit_pipeline/SELECTED_PAIR_SAMPLING.md}
\item \path{tests/test_epit_selected_sampler.py}
\item \path{tests/test_epit_selected_pairs.py}
\item \path{tests/test_epit_selected_training_command.py}
\end{itemize}
The earlier diagnostic counts come from
\path{corrosion_datasets/analysis/epit_pipeline/formula_mixture_real_features/cv5_comparison/}.
The new pair counts come from the separate
\path{selected_pairs_all760_generic10k/} preparation outputs.

\end{document}
